% House-style research paper template (research-paper-writer-pro v2).
% A4, 11pt, single column, Palatino text and math, author–year APA references,
% numbered propositions and equations, booktabs tables, figures with captions below.
% Build: latexmk -pdf -interaction=nonstopmode paper.tex   (bibtex runs automatically)

\documentclass[11pt,a4paper]{article}

% ---- Typography ----------------------------------------------------------
\usepackage[T1]{fontenc}
\usepackage[utf8]{inputenc}
\usepackage{amsmath,amsthm}          % before newpxmath, which supplies the AMS symbol set itself
\usepackage{newpxtext,newpxmath}     % Palatino text + matching math
\usepackage[margin=1in]{geometry}
\usepackage{microtype}
\usepackage{parskip}                 % space between paragraphs, no first-line indent
\usepackage{booktabs}
\usepackage{array}
\newcolumntype{L}[1]{>{\raggedright\arraybackslash}p{#1}}  % ragged text columns for tables
\usepackage{graphicx}
\usepackage[font=small,labelfont=bf,labelsep=colon]{caption}
\usepackage{xcolor}
\usepackage{csquotes}
\usepackage{enumitem}

% ---- References: APA author–year in text, alphabetical list (BibTeX) -----
% apacite + BibTeX builds on any TeX Live. biblatex-apa (style=apa, backend=biber) is the
% exact-APA-7 alternative when a working biber exists; swap this block, \citep->\parencite,
% \citet->\textcite, and \bibliography->\printbibliography.
\usepackage[natbibapa]{apacite}
\bibliographystyle{apacite}
\AtBeginDocument{\setlength{\bibitemsep}{0.45\baselineskip}}   % air between entries

% ---- Running head and page numbers ---------------------------------------
\newcommand{\shorttitle}{Short running title}
\newcommand{\shortauthor}{A. Author}
\usepackage{fancyhdr}
\pagestyle{fancy}
\fancyhf{}
\lhead{\small\shorttitle}
\rhead{\small\shortauthor}
\cfoot{\thepage}
\renewcommand{\headrulewidth}{0.4pt}
\setlength{\headheight}{14pt}

% ---- Links: one navy for internal, citation and URL links ----------------
\usepackage[hidelinks]{hyperref}
\definecolor{linknavy}{RGB}{0,0,112}
\hypersetup{colorlinks=true,linkcolor=linknavy,citecolor=linknavy,urlcolor=linknavy}
\AtBeginDocument{%
  \urlstyle{same}%                                   URLs in the text face, not monospace
  \renewcommand{\doiprefix}{}%                      APA 7 prints the DOI as a bare link
  \renewcommand{\doi}[1]{\href{https://doi.org/#1}{https://doi.org/#1}}}

% ---- Numbered environments (bold label, italic body for results) ---------
\theoremstyle{plain}
\newtheorem{proposition}{Proposition}
\newtheorem{corollary}{Corollary}
\newtheorem{lemma}{Lemma}
\newtheorem{theorem}{Theorem}
\theoremstyle{definition}
\newtheorem{definition}{Definition}
\newtheorem{assumption}{Assumption}
\theoremstyle{remark}
\newtheorem{remark}{Remark}

% ---- Title block ---------------------------------------------------------
\title{\bfseries Main Title of the Paper\\[0.5ex]
  {\large\mdseries A Subtitle Stating the Contribution or Framing}}
\author{Author Name\\[0.3ex]
  {\normalsize Affiliation, City, Country}\\[0.3ex]
  {\normalsize\href{mailto:author@example.org}{author@example.org}}}
\date{\normalsize \today}

\begin{document}
\maketitle
\thispagestyle{plain}   % page number, no running head, on the title page

\begin{abstract}
\noindent One paragraph of 200--300 words. State the problem, the approach, each
result as a conditional statement, the worked application, and what the
contribution is and is not. No citations, no undefined symbols.
\end{abstract}

\noindent\textbf{Keywords:} first keyword; second keyword; third keyword; fourth keyword

\section{Introduction}

Open with the question in one or two plain sentences. Situate it in the literature
with author--year citations in parentheses \citep{jervis1978, harsanyi1988}, or
in the sentence: \citet{sastry2024} argue that compute is a governance lever.
State the narrower question this paper answers, in italics if it is a single
sentence. Then list the contributions (``The framework makes three contributions.
First, \ldots'') and close with what the contributions are and are not.

\section{Research design and theoretical foundations}

\subsection{Research design}

Say what kind of paper this is and why the theories were selected. Give the
analytical procedure as numbered steps in prose (``The analytical procedure has
four steps. The first \ldots'').

\subsection{Foundations}

Introduce notation in displayed, numbered equations and refer to them by number:
\begin{equation}
  U_i = \pi_i - \alpha\,\pi_j, \qquad \alpha = \frac{k}{1+k} \in [0,1).
  \label{eq:positional}
\end{equation}
Equation~\eqref{eq:positional} is then cited by number, never by ``the equation above''.

\section{Framework and units of analysis}

Define every object and level of analysis before any result. A table that maps
mechanisms to their inputs, implications and open empirical questions belongs
here (Table~\ref{tab:mechanisms}).

\begin{table}[t]
  \centering
  \caption{Mechanisms, assumptions, and inferential boundaries.}
  \label{tab:mechanisms}
  \small
  \begin{tabular}{@{}L{2.4cm}L{3.1cm}L{3.9cm}L{3.9cm}@{}}
    \toprule
    Mechanism & Formal inputs & Conditional implication & Additional empirical question \\
    \midrule
    First mechanism & Inputs named as symbols & What follows, stated as a bound & What would have to be measured \\
    Second mechanism & \ldots & \ldots & \ldots \\
    \bottomrule
  \end{tabular}
\end{table}

\section{Formal analysis}

\subsection{First result}

\begin{proposition}[Descriptive name]
\label{prop:first}
State the result with every quantifier and assumption inside the statement. Bound
its scope in the last sentence (``This is a bound on \ldots, not an estimate of
\ldots'').
\end{proposition}

Follow each result with plain prose on what it does and does not establish.
Illustrate with a figure whose numbers are declared illustrative
(Figure~\ref{fig:first}).

\begin{figure}[t]
  \centering
  \includegraphics[width=0.9\textwidth]{figures/fig1.pdf}
  \caption{Illustrative upper bounds on the quantity of interest with $d = 150$.
  These are bounds, not estimates.}
  \label{fig:first}
\end{figure}

\begin{corollary}[A joint condition]
\label{cor:joint}
A corollary reads the same way: statement first, scope last.
\end{corollary}

\section{Interpretation}

Connect the results to each other by number (``Corollary~\ref{cor:joint} identifies
\ldots; Proposition~\ref{prop:first} identifies a different \ldots'') and say which
mistaken substitutions the separation prevents.

\section{Worked application}

\subsection{Define the activity before assigning parameters}

\subsection{A numerical diagnostic, not a calibration}

All numbers assumed; say so in the table caption.

\section{Empirical development and falsifiable implications}

\subsection{Separate the tests}

\subsection{Identification and the observation model}

\subsection{Discriminating implications}

\section{Discussion and conclusion}

Restate what the framework distinguishes, what it does not claim, and the
empirical priorities in order.

\section*{Data, code, and AI assistance}

Name the scripts that produced every figure and table, where they live, and what
AI assistance was used for.

\appendix
\section{Proofs and supplementary derivations}

\subsection{Proof of Proposition~\ref{prop:first}}

\begin{proof}
Proofs go here, one subsection per result. The box closes each proof.
\end{proof}

\section*{Numerical illustration and reproducibility}

Parameters and seeds for every figure, so a reader can regenerate them.

\bibliography{references}

\end{document}
